% !TeX root = main.tex
\chapter{Introduction}

\quad Machine translation evolved a lot in the last years, 
but my question was if the subtask of code translation
evolved at the same level. Translating programs from one
language to another can support code migration, reuse, and maintenance, but the
task is harder than ordinary text translation. A correct program translation
must preserve executable behavior: the generated code has to compile, run, and
produce the same outputs as the source program.

This thesis studies Java-Python translation on the AVATAR dataset, which
contains aligned competitive-programming solutions and execution test cases. The
task is evaluated in both directions, Java-to-Python and Python-to-Java. The
main challenge is that lexical similarity is not enough: a translation can
obtain reasonable BLEU or CodeBLEU scores while still failing to preserve the
required algorithmic behavior. For this reason, the evaluation also includes
compilation or syntax checking, Execution Accuracy, and Computational Accuracy.

This work compares one modern model for each required category: zero-shot Qwen
7B, encoder-only UniXcoder with a learned Transformer decoder, decoder-only Qwen
0.5B with QLoRA, and Seq2Seq CodeT5+ 220M. The strongest final system is the
zero-shot Qwen model, reaching 59.68
Computational Accuracy for Java-to-Python and 61.80 for Python-to-Java. Among
the trained systems, the Seq2Seq CodeT5+ model gives the best functional
result, reaching 47.82 Computational Accuracy for Java-to-Python and 41.29 for
Python-to-Java after two epochs of fine-tuning per direction.

Compared with the original AVATAR baselines, this work focuses on modern
pretrained code models. The comparison shows the practical effect of code pretraining,
instruction tuning, and parameter-efficient fine-tuning under the same benchmark
and metric protocol. It also shows why execution-based evaluation is necessary:
modern models can generate code that is textually close to the reference while
still failing to preserve the program behavior.
